\begin{table*}[t]\centering\small\setlength{\tabcolsep}{5pt}
\begin{adjustbox}{max width=\textwidth}\begin{tabular}{lccccccccc}
\toprule
 & \multicolumn{3}{c}{PUBHEALTH} & \multicolumn{3}{c}{CLIMATE-FEVER} & \multicolumn{3}{c}{Fin-Fact} \\
\cmidrule(lr){2-4} \cmidrule(lr){5-7} \cmidrule(lr){8-10}
Detector & Base & Inj. & Decoy & Base & Inj. & Decoy & Base & Inj. & Decoy \\
\midrule
Claim retrieval + LLM \citep{pikuliak2023multiclaim} & 63.3 & 61.3 & 56.0 & 45.9 & 47.8 & 45.2 & 65.8 & 64.0 & 57.6 \\
XLM-R, adv.\ training \citep{conneau2020xlmr} & 72.1 & 66.4 & 66.8 & 71.8 & 66.9 & 58.6 & 74.6 & 67.7 & 67.4 \\
\method{} (ours) & 75.9 & 72.8 & 67.2 & 68.1 & 65.3 & 56.8 & 76.6 & 71.4 & 62.6 \\
\method{} + \air{} + aug. (ours) & 77.3 & 70.9 & 65.5 & 66.2 & 54.4 & 48.4 & 79.3 & 70.4 & 62.3 \\
\bottomrule
\end{tabular}\end{adjustbox}
\caption{Adaptive attacks. Macro-F1 (\%) averaged over the summarisation and MCQ$\rightarrow$JSON attacks, without (Base) and with an injected instruction to the claim extractor (Inj.) or three true decoy claims around the content (Decoy). \method{} + \air{} + aug.: the Setting~B variant of \method{}, with the attack-invariant retriever (\air{}) and a verifier also trained on the transformed (augmented) training claims.}
\label{tab:adaptive}
\end{table*}
